\textbf{Seats}
\begin{enumerate}
    \item 5 mins: Intro slides \par
    \textbf{Motivation}
    \item 18 mins: Finish partial fractions where we left off \par
    \textbf{Check piazza}
    \item 24 mins: Finish Section 2.1 \par
    \textbf{Check piazza}
    \item 36 mins: Finish Example 2.2.1 \par
    \textbf{Check piazza}
    \item 45 mins: Finish Example 2.2.2 \par
    \textbf{Check piazza}
    \item 50 mins: Wrapped up
\end{enumerate}
\textbf{Seats}

\vspace{10mm}
\textbf{Example 1.3.5}
Factorise the denominator $3(x+1)^2/(x^2+x-2) = 3(x+1)^2/[(x - 1)(x + 2)]$, then
\begin{align*}
    \frac{3(x+1)^2}{(x - 1)(x + 2)} &= A + \frac{B}{x-1} + \frac{C}{x+2}\\
    \therefore 3(x+1)^2 = A(x-1)(x+2) + B(x+2) + C(x-1)
\end{align*}
Set $x=1 \quad \rightarrow \quad 3B=12 \quad \rightarrow \quad B=4$\\
Set $x=-2 \quad \rightarrow \quad -3C=3 \quad \rightarrow \quad C=-1$\\
Comparing coefficients of $x^2$: $A=3$\\
Therefore
\begin{align*}
    \frac{3(x+1)^2}{x^2+x-2} = 3 + \frac{4}{x-1} - \frac{1}{x+2}
\end{align*}

\vspace{10mm}
\textbf{Example 1.3.6}
\begin{itemize}
    \item Order of numerator is 6 and denominator is 4. So need a quadratic.
    \item Repeated root so need a $(x-2)$ and $(x-2)^2$ denominator.
    \item Irreducible quadratic so need a linear term in numerator of this.
\end{itemize}
Therefore
\begin{align*}
    Ax^2 + Bx + C + \frac{D}{x-2} + \frac{E}{(x-2)^2} + \frac{Fx+G}{x^2-6x+50}.
\end{align*}

\vspace{10mm}
\textbf{Example 2.2.1}
\begin{align*}
f(x) &= 2x^3 - 9x^2 + 5 \\
f'(x) &= 6x^2 -18x\\
&= 6x(x-3)
\end{align*}
So, stationary points $x=0$ and $x=3$. 

To use the first derivative test, let's choose $h=1$ (as this is a small enough step to not cross another stationary point, but also makes the arithmetic simple).
\begin{itemize}
    \item Around $x = 0$: $f'(-1) = 24 > 0$ and $f'(1)=-12<0$. Therefore, /$\overline{\phantom{x}}$\textbackslash, and so the point is a maximum.
    \item Around $x = 3$: $f'(2) = -12 < 0$ and $f'(4)=24>0$. Therefore, \textbackslash$\underline{\phantom{x}}$/, and so the point is a minimum.
\end{itemize}
Second derivative test: $f''(x) = 12x - 18$. 
\begin{itemize}
    \item $f''(0) = -18 < 0$, therefore maximum\\
    \item $f''(3) = 18 > 0$, therefore minimum
\end{itemize}

\vspace{10mm}
\textbf{Example 2.2.2}
First, find the stationary points $p'(x) = 0$.
\begin{align*}
    p(x) &= 2x^3-3x^2-36x+13\\
    p'(x) &= 6x^2-6x-36\\
    &= 6(x-3)(x+2).
\end{align*}
So, stationary points at $x = \{-2,3\}$.

The function is continuous so we just need to check
\begin{itemize}
    \item The left-hand limit: $x=-4 \quad \rightarrow \quad p(-4) = 2(-64) - 3(16) - 36(-4) + 13 = -19$
    \item The right-hand limit: $x=6 \quad \rightarrow \quad p(6) = 2(216)-3(36)-36(6)+13 = 121$
    \item The first stationary point $x=-2 \quad \rightarrow \quad p(-2) = 2(-8) - 3(4) - 36(-2) + 13 = 57$
    \item The second stationary point $x=3 \quad \rightarrow \quad p(3) = (27)-3(9)-36(3)+13=-68$
\end{itemize}
So, 
\begin{itemize}
    \item Minimum is $p=-68$ at $x=3$.
    \item Maximum is $p = 121$ at $x=6$.
\end{itemize}
Therefore, the minimum of $p(x)$ over the interval $-4 \leq x \leq 6$ occurs at the stationary point $x=3$, where the function takes the value $p(3) = -68$. The maximum occurs at the upper limit of the interval $x=6$, where $p(6) = 121$.



